% GENERATED FILE. Edit canonical lessons, then run npm run generate:books.
% canonical-source-hash: fnv1a64:0738adc2bd321019
% canonical-lessons: HI-C216-mugphali, HI-C216-muli, HI-C216-baigan, HI-C216-bhindi, HI-C216-papita

\chapter{A peanut, A radish, An aubergine, Okra, A papaya}
\label{ch:hi-a-peanut-a-radish-an-aubergine-okra-a-papaya}
\hlchaptermodality{216}

\begin{chapteropening}
\textbf{By the end of this chapter:} \emph{I can say a peanut, a radish, an aubergine, okra and a papaya.}

Complete the last lesson of chapter 216: i can say a peanut, a radish, an aubergine, okra and a papaya.
\end{chapteropening}

\section[\texorpdfstring{\hi{मूँगफली} (mū̃gphalī)}{मूँगफली (mū̃gphalī)}]{\hi{मूँगफली} (mū̃gphalī) — a peanut}

\label{lesson:HI-C216-mugphali}



\begin{quote}
Before the new one: say the Hindi for brave, then the Hindi for kind.
\end{quote}

\subsection*{You'll want to know: \hi{मूँगफली}}
\textbf{\hi{मूँगफली}} — \emph{mū̃gphalī} — \textquotedblleft{}a peanut\textquotedblright{}.

\begin{grammarlens}[title={What you've built}]
One more word for this chapter.
\end{grammarlens}

\subsection*{Your turn}
\begin{itemize}
\raggedright
  \item \emph{Say it:} \emph{mū̃gphalī}
  \item \emph{Say it:} \emph{mū̃gphalī}, once more
  \item \emph{Say it:} say \emph{mū̃gphalī}
  \item \emph{Recall:} say the Hindi for brave, then the Hindi for kind, then say \emph{mū̃gphalī} again
\end{itemize}

\begin{tcolorbox}[breakable,colback=teal!4,colframe=teal!35!black,title={Before you move on}]
What does \hi{मूँगफली} mean? (\textquotedblleft{}A peanut\textquotedblright{}.) Say it once more.
\end{tcolorbox}
\section[\texorpdfstring{\hi{मूली} (mūlī)}{मूली (mūlī)}]{\hi{मूली} (mūlī) — a radish}

\label{lesson:HI-C216-muli}



\begin{quote}
Before the new one: say the Hindi for kind, then the Hindi for a peanut.
\end{quote}

\subsection*{You'll want to know: \hi{मूली}}
\textbf{\hi{मूली}} — \emph{mūlī} — \textquotedblleft{}a radish\textquotedblright{}.

\begin{grammarlens}[title={What you've built}]
One more word for this chapter.
\end{grammarlens}

\subsection*{Your turn}
\begin{itemize}
\raggedright
  \item \emph{Say it:} \emph{mūlī}
  \item \emph{Say it:} \emph{mūlī}, once more
  \item \emph{Say it:} say \emph{mūlī}
  \item \emph{Recall:} say the Hindi for kind, then the Hindi for a peanut, then say \emph{mūlī} again
\end{itemize}

\begin{tcolorbox}[breakable,colback=teal!4,colframe=teal!35!black,title={Before you move on}]
What does \hi{मूली} mean? (\textquotedblleft{}A radish\textquotedblright{}.) Say it once more.
\end{tcolorbox}
\section[\texorpdfstring{\hi{बैंगन} (bãigan)}{बैंगन (bãigan)}]{\hi{बैंगन} (bãigan) — an aubergine}

\label{lesson:HI-C216-baigan}



\begin{quote}
Before the new one: say the Hindi for a peanut, then the Hindi for a radish.
\end{quote}

\subsection*{You'll want to know: \hi{बैंगन}}
\textbf{\hi{बैंगन}} — \emph{bãigan} — \textquotedblleft{}an aubergine\textquotedblright{}.

\begin{grammarlens}[title={What you've built}]
One more word for this chapter.
\end{grammarlens}

\subsection*{Your turn}
\begin{itemize}
\raggedright
  \item \emph{Say it:} \emph{bãigan}
  \item \emph{Say it:} \emph{bãigan}, once more
  \item \emph{Say it:} say \emph{bãigan}
  \item \emph{Recall:} say the Hindi for a peanut, then the Hindi for a radish, then say \emph{bãigan} again
\end{itemize}

\begin{tcolorbox}[breakable,colback=teal!4,colframe=teal!35!black,title={Before you move on}]
What does \hi{बैंगन} mean? (\textquotedblleft{}An aubergine\textquotedblright{}.) Say it once more.
\end{tcolorbox}
\section[\texorpdfstring{\hi{भिंडी} (bhinḍī)}{भिंडी (bhinḍī)}]{\hi{भिंडी} (bhinḍī) — okra}

\label{lesson:HI-C216-bhindi}



\begin{quote}
Before the new one: say the Hindi for a radish, then the Hindi for an aubergine.
\end{quote}

\subsection*{You'll want to know: \hi{भिंडी}}
\textbf{\hi{भिंडी}} — \emph{bhinḍī} — \textquotedblleft{}okra\textquotedblright{}.

\begin{grammarlens}[title={What you've built}]
One more word for this chapter.
\end{grammarlens}

\subsection*{Your turn}
\begin{itemize}
\raggedright
  \item \emph{Say it:} \emph{bhinḍī}
  \item \emph{Say it:} \emph{bhinḍī}, once more
  \item \emph{Say it:} say \emph{bhinḍī}
  \item \emph{Recall:} say the Hindi for a radish, then the Hindi for an aubergine, then say \emph{bhinḍī} again
\end{itemize}

\begin{tcolorbox}[breakable,colback=teal!4,colframe=teal!35!black,title={Before you move on}]
What does \hi{भिंडी} mean? (\textquotedblleft{}Okra\textquotedblright{}.) Say it once more.
\end{tcolorbox}
\section[\texorpdfstring{\hi{पपीता} (papītā)}{पपीता (papītā)}]{\hi{पपीता} (papītā) — a papaya}

\label{lesson:HI-C216-papita}



\begin{quote}
Before the new one: say the Hindi for an aubergine, then the Hindi for okra.
\end{quote}

\subsection*{You'll want to know: \hi{पपीता}}
\textbf{\hi{पपीता}} — \emph{papītā} — \textquotedblleft{}a papaya\textquotedblright{}.

\begin{grammarlens}[title={What you've built}]
One more word for this chapter.
\end{grammarlens}

\subsection*{Your turn}
\begin{itemize}
\raggedright
  \item \emph{Say it:} \emph{papītā}
  \item \emph{Say it:} \emph{papītā}, once more
  \item \emph{Say it:} say \emph{papītā}
  \item \emph{Recall:} say the Hindi for an aubergine, then the Hindi for okra, then say \emph{papītā} again
\end{itemize}

\begin{tcolorbox}[breakable,colback=teal!4,colframe=teal!35!black,title={Before you move on}]
What does \hi{पपीता} mean? (\textquotedblleft{}A papaya\textquotedblright{}.) Say it once more.
\end{tcolorbox}
